\bookchapter{Глава 9: Мобильный офис и обучение}

\chapterintro{Раньше я тоже считал, что работать на телефоне — это «на скорую руку», пока однажды в командировке не забыл ноутбук и был вынужден править тендерную документацию и подписывать договоры исключительно на телефоне. Честно говоря, узкое место мобильной работы часто не в самом телефоне, а в том, что вы не знакомы с его реальными возможностями.}

\booksection{Обработка документов: редактирование и творчество на ладони}

Редактировать документы на экране телефона — ощущения действительно иные, чем на компьютере. Главная проблема — маленький экран, толстые пальцы, вёрстка легко сбивается. Но несколько приёмов могут значительно улучшить опыт:

\begin{itemize}
  \item \textbf{WPS Office}: лучшая локализация, полная совместимость с форматами Microsoft (docx/xlsx/pptx), богатый встроенный инструментарий PDF. Мне особенно нравится функция «Голосовые заметки» — когда неудобно печатать при просмотре контракта, просто надиктуйте замечания в телефон, и они автоматически превратятся в текстовые аннотации в документе. Кроме того, в WPS встроено огромное количество бесплатных шаблонов, которые экономят уйму времени при создании резюме и презентаций.
  \item \textbf{Microsoft 365}: если в вашей компании используется лицензионный Office, это лучший выбор. Совместимость форматов с десктопной версией почти идеальная — в отличие от WPS, где сложная вёрстка может слегка смещаться. Глубокая интеграция с OneDrive: документ, который вы редактировали на компьютере, открывается на телефоне в последней версии.
  \item \textbf{Shimo Docs / Tencent Docs}: это представители «онлайн-коллаборации». Подходят для одновременного редактирования одного документа несколькими людьми — изменения каждого отображаются в реальном времени. Я использую Shimo Docs с друзьями для совместного планирования путешествий: двое одновременно заполняют таблицу — намного эффективнее, чем одному.
\end{itemize}

\textbf{Практичный рабочий процесс}: моя привычка сейчас — важные файлы хранить в OneDrive, повседневное редактирование в WPS, командную работу в Shimo Docs. С тех пор как файлы переехали в облако, я больше ни разу не попадал в неловкую ситуацию «забыл флешку».

\booksection{Приложения для заметок: от быстрых записей к системе знаний}

Многие, начиная вести заметки, сразу пытаются построить сложную систему — и бросают через три дня. Мой совет: от простого к сложному — сначала выработайте привычку «записывать», потом думайте об «управлении».

\begin{itemize}
  \item \textbf{Быстрые записи — встроенных достаточно}: Заметки iOS запускаются молниеносно, поддерживают сканирование документов, рукописный ввод и запись голоса. Фирменные блокноты Android тоже неплохи — можно быстро запустить даже с экрана блокировки. Эти инструменты хороши для мимолётных идей, списков покупок, ключевых моментов встреч. Главное — быстрый запуск: подумал — записал, не дожидаясь загрузки приложения.
  \item \textbf{Отложенное чтение — сбор информации}: функция «Плавающее окно» в WeChat — на самом деле скрытый шедевр. Увидели хорошую статью, но нет времени читать — нажмите «Плавающее окно» в правом верхнем углу, и она свернётся в маленький кружок на краю экрана, готовый к прочтению в любой момент. Гораздо удобнее, чем пересылать в «Помощник по передаче файлов». Кроме того, инструменты вроде Cubox или Pocket позволяют собирать интересный контент со всех платформ.
  \item \textbf{Построение системы — Notion и Feishu}: когда разрозненные заметки становится трудно искать, пора переходить на Notion. Эта штука невероятно гибкая — можно создать «Базу данных книжных заметок» с автоматической организацией по названию, автору, оценке; или построить «Доску проекта» для визуализации прогресса задач. «Многомерная таблица» в Feishu похожа, и в Китае работает быстрее. Я управляю всеми своими материалами для письма, учебными заметками и прогрессом проектов в Notion уже более трёх лет.
\end{itemize}

\textbf{Практичная стратегия заметок}: моя привычка — «тройной комплект»: мимолётные идеи — в системный блокнот (мгновенный запуск), статьи для чтения — в плавающее окно WeChat (под рукой), знания для долгосрочного накопления — в Notion (систематизация). Не пытайтесь с самого начала запихнуть всё в один инструмент — так легче бросить.

\booksection{Сканирование и OCR: перенос бумажного мира в телефон}

Если правильно использовать сканирование на телефоне, можно избавиться от множества походов в копировальный центр. Когда я оформлял визу и нужна была куча электронных копий документов — всё сделал на телефоне.

\begin{itemize}
  \item \textbf{Пользователям iPhone}: встроенная функция сканирования в Заметках iOS уже очень хороша. Откройте Заметки нажмите значок камеры выберите «Сканировать документы», наведите на документ — он автоматически распознает границы и скорректирует перспективу. После съёмки можно настроить цвет: режим «Чёрно-белый» даёт результат, наиболее близкий к сканеру, «Серый» подходит для материалов с картинками.
  \item \textbf{Пользователям Android}: многие брендовые камеры имеют встроенный режим «Сканирование документов». Если не нашли — скачайте Microsoft Lens из магазина приложений (бесплатно, с интеграцией в OneDrive).
  \item \textbf{OCR-распознавание текста}: если нужно извлечь текст из отсканированного документа для редактирования, нужен OCR. Встроенная функция WeChat «Извлечь текст» — самая простая: долгое нажатие на картинку, и точность распознавания неплохая. Для больших объёмов можно использовать приложение Baimiao или CamScanner. Для китайского текста OCR-движок Baidu работает лучше; функция WPS «Картинка в текст» тоже очень точна.
  \item \textbf{Советы по сканированию}: снимайте при хорошем освещении, лучше с настольной лампой сбоку, чтобы тень от руки не падала на документ. Если уголок загнут — сначала прижмите его рукой. Многостраничные документы можно сканировать последовательно и объединить в один PDF, не сохраняя каждую страницу отдельно.
\end{itemize}

\booksection{Как выбрать облачное хранилище: нет лучшего, есть подходящее именно вам}

Я долго размышлял над выбором облачного хранилища и в итоге понял: всё зависит от вашей «основной экосистемы».

\begin{itemize}
  \item \textbf{Экосистема Apple iCloud}: бесшовная синхронизация настолько незаметна, что вы её не чувствуете. Сфотографировали на телефон — через несколько минут фото уже на компьютере. И защита приватности у Apple самая строгая. Но очевидный минус: делиться с пользователями Android или Windows неудобно, бесплатно всего 5 ГБ, расширение платное.
  \item \textbf{Основной Windows OneDrive}: если у вас Windows-компьютер + Office 365, OneDrive — идеальная пара. Сохраняете документ в Word — он автоматически в облаке. 1 ТБ пространства с подпиской Office — очень выгодно. Автозагрузка фото в мобильном приложении тоже отлично работает.
  \item \textbf{Обмен файлами в Китае Baidu Wangpan}: честно говоря, ограничение скорости скачивания у Baidu Wangpan постоянно критикуют, но его экосистема обмена действительно вне конкуренции. Огромное количество учебных материалов и медиаресурсов в китайском интернете распространяется именно через Baidu Wangpan. Если вы лишь изредка скачиваете чужие файлы, можно обойтись без платной подписки. Если же вы активно коллекционируете ресурсы — оформляйте подписку, сэкономленное время стоит гораздо дороже.
  \item \textbf{Alibaba Cloud Drive / Quark Drive}: восходящие звёзды последних двух лет. Alibaba Cloud Drive не ограничивает скорость загрузки и скачивания (пока), а Quark Drive глубоко интегрирован с браузером Quark — можно сохранять материалы прямо в облако во время поиска, очень удобно.
\end{itemize}

\textbf{Моя комбинированная схема}: iCloud — для личных фото и повседневных файлов (приоритет приватности), Baidu Wangpan — для ресурсов и приёма общих файлов (приоритет экосистемы), OneDrive — для рабочих документов (приоритет интеграции с Office). Три облачных хранилища, каждое выполняет свою роль, не мешая друг другу.

\booksection{Реальный сценарий мобильной работы}

Расскажу о своём реальном опыте. Однажды в командировке клиент внезапно потребовал изменить условия договора. Мой рабочий процесс выглядел так:

\begin{enumerate}
  \item В отеле получил от клиента правки в WeChat (текстовое сообщение + фото с рукописными пометками).
  \item Открыл OneDrive, нашёл оригинал договора, открыл в WPS для редактирования и постатейно внёс изменения, сверяясь с сообщениями клиента в WeChat.
  \item После правок экспортировал через функцию WPS «Экспорт в PDF» и отправил в групповой чат WeChat клиенту на утверждение.
  \item Клиент прислал второй раунд замечаний — я использовал инструмент «Кисть» в WPS, чтобы прямо в PDF обвести места изменений, и добавил голосовые заметки с пояснением логики правок.
  \item После подтверждения подписал онлайн через функцию «Электронная подпись» в WPS и отправил контрагенту для встречной подписи.
  \item Наконец, сохранил подписанный PDF-договор в Baidu Wangpan и iCloud (по копии в каждом), а также сделал запись в рабочем журнале Notion.
\end{enumerate}

Весь процесс от начала до конца занял около сорока минут, и я ни разу не открыл ноутбук. Честно говоря, три года назад я бы не поверил, что телефон на такое способен, но сегодня это реально. Ключевой момент — заранее освоить эти инструменты, а не судорожно скачивать и регистрироваться, когда они срочно понадобятся.

\begin{tcbpractice}{Попробуйте: постройте свой рабочий процесс мобильного офиса}
\begin{enumerate}
  \item Выберите основное приложение для документов (WPS или Microsoft 365), войдите и подключите облачное хранилище.
  \item Откройте Заметки или Блокнот, отсканируйте документ и экспортируйте как PDF, отправив себе.
  \item Выберите приложение облачного хранилища и настройте автоматическое резервное копирование фото (выберите «Только по Wi-Fi» для экономии трафика).
  \item Создайте новую страницу в Notion или Feishu и попробуйте построить простой шаблон «Рабочего журнала».
\end{enumerate}

\end{tcbpractice}
